\documentclass[../DoAn.tex]{subfiles}
\begin{document}

% =========================================================
\section{Mở đầu chương}
% =========================================================

Các chương trước đã trình bày kiến trúc, thiết kế và các đóng góp kỹ thuật của
hệ thống trợ lý học vụ. Chương này trình bày kết quả thực nghiệm đánh giá hệ
thống theo hướng đầu cuối (end-to-end): mỗi câu hỏi được chạy qua toàn bộ luồng
xử lý thực tế của pipeline --- từ định tuyến, viết lại truy vấn, truy hồi lai,
xếp hạng lại, lọc hiệu lực văn bản, cho đến sinh câu trả lời và tự đánh giá
chất lượng đầu ra.

Nội dung chương gồm: thiết lập thực nghiệm, định nghĩa các chỉ số đánh giá, kết
quả tổng thể khi so sánh các cấu hình truy hồi, đánh giá tách thành phần
(ablation study), phân tích kết quả, hạn chế còn tồn tại và hướng phát triển
tiếp theo.

% =========================================================
\section{Thiết lập thực nghiệm}
% =========================================================

\subsection{Hệ thống được đánh giá}

Hệ thống được đánh giá trong chương này là toàn bộ pipeline trợ lý học vụ với
cấu hình gần với môi trường vận hành thực tế: bật tác nhân tự động, bật cơ chế
truy hồi bổ sung bằng tài liệu giả định (HyDE), bật bộ lọc hiệu lực văn bản,
ngưỡng điểm xếp hạng lại ở mức 0,0, số đoạn tối thiểu sau xếp hạng lại là 3,
và lấy tối đa 7 đoạn cho mỗi lượt truy hồi. Thiết lập ngưỡng 0,0 được chọn có
chủ đích để giữ lại tối đa tài liệu sau bước xếp hạng lại, qua đó đánh giá năng
lực xếp hạng của mô hình thay vì loại bỏ ứng viên quá sớm. Khác biệt duy nhất
so với hệ thống thực là tắt bộ nhớ đệm để mỗi truy vấn luôn sinh câu trả lời
mới hoàn toàn. Nhờ đó, số liệu phản ánh hành vi đầu cuối mà người dùng nhận
được thay vì đo kết quả từ bộ nhớ đệm.

Bộ đánh giá gọi trực tiếp vào điểm vào chính của pipeline chat --- cùng điểm vào
với API thực tế --- và ghi lại tập đoạn truy hồi được (để tính chỉ số truy hồi),
câu trả lời sinh ra, cùng các nhãn đánh giá đầu cuối do mô hình đánh giá và bộ
so khớp đáp án vàng gán.

Lưu ý về cơ chế HyDE: Vì ngưỡng xếp hạng lại được đặt ở mức 0,0, điều kiện
kích hoạt HyDE (toàn bộ ứng viên có điểm logit âm) trở thành một cơ chế bảo vệ
dự phòng khắt khe cho các câu hỏi nhiễu. Trong các lượt thực nghiệm được dùng
cho báo cáo này, chất lượng truy hồi sơ bộ đủ tốt để tài liệu trả về vượt qua
ngưỡng 0,0; do đó các chỉ số đo lường chủ yếu phản ánh năng lực của luồng truy
hồi và xếp hạng lại cốt lõi.

\subsection{Bộ dữ liệu thực nghiệm}

Thực nghiệm được tiến hành trên tập dữ liệu đánh giá bao quát cả bốn kho tri
thức: chương trình đào tạo, quy định học vụ, kế hoạch học kỳ và sổ tay sinh
viên. Các câu hỏi được thiết kế để kiểm tra khả năng trả lời chính xác các
thông tin học vụ, từ các câu hỏi tra cứu thông tin đơn giản đến các truy vấn
đòi hỏi tổng hợp từ nhiều nguồn.

Tại thời điểm tổng hợp báo cáo, bộ \texttt{rag\_evaluation\_dataset} đang được
đánh giá lại ở cấu hình RRF cân bằng trọng số. Vì vậy, chương này tạm thời sử
dụng 15 bộ dữ liệu còn lại với tổng cộng 470 truy vấn. Cách chọn này bảo đảm
tất cả cấu hình được so sánh trên cùng một phạm vi dữ liệu, tránh sai lệch do
thiếu kết quả ở một cấu hình riêng lẻ.

\subsection{Các cấu hình được so sánh}

Như đã trình bày ở Chương~2, hệ thống hỗ trợ hợp nhất kết quả truy hồi từ hai
nguồn: tìm kiếm vector và tìm kiếm từ khóa BM25. Chương này so sánh ba nhóm
cấu hình chính. Nhóm thứ nhất đánh giá chiến lược hợp nhất tuyến tính và RRF
với cùng trọng số mặc định 0,8 cho vector và 0,2 cho từ khóa. Nhóm thứ hai khảo
sát RRF với trọng số cân bằng 0,5/0,5, trong đó kết quả đã được đánh giá lại
bằng Gemini. Nhóm thứ ba là các thí nghiệm tách thành phần gồm vector-only,
keyword-only và no-rerank.

\begin{table}[H]
\centering
\small
\begin{tabular}{|p{2.4cm}|p{5.6cm}|p{5.6cm}|}
\hline
& \textbf{Hợp nhất tuyến tính} & \textbf{Hợp nhất theo thứ hạng (RRF)} \\
\hline
\textbf{Nguyên lý} & Chuẩn hóa mỗi tập về $[0,1]$ theo giá trị cao nhất, rồi
cộng có trọng số: $0{,}8 \cdot \widehat{S_v} + 0{,}2 \cdot \widehat{S_k}$ &
Hợp nhất theo thứ hạng: $0{,}8 \cdot \frac{1}{k + r_v} + 0{,}2 \cdot
\frac{1}{k + r_k}$ với $k = 10$ \\
\hline
\textbf{Phụ thuộc vào} & Độ lớn điểm tuyệt đối & Chỉ thứ hạng, bỏ qua độ lớn \\
\hline
\textbf{Điểm yếu} & Nhạy khi thang điểm cosine và BM25 lệch nhau & Nén khoảng
cách giữa các hạng liền kề \\
\hline
\end{tabular}
\caption{So sánh hai cơ chế hợp nhất kết quả truy hồi}
\label{tab:ch6-fusion}
\end{table}

Cấu hình RRF 0,8/0,2 là cấu hình vận hành chính được đánh giá trong chương
này. Cấu hình RRF 0,5/0,5 được xem là thí nghiệm bổ sung nhằm khảo sát độ nhạy
của hệ thống khi tăng tỷ trọng tín hiệu từ khóa. Các cấu hình vector-only,
keyword-only và no-rerank được dùng để làm rõ vai trò của từng thành phần trong
pipeline.

% =========================================================
\section{Định nghĩa chỉ số đánh giá}
% =========================================================

\subsection{Chỉ số truy hồi}

Các chỉ số truy hồi đối chiếu tập định danh đoạn truy hồi được với tập bằng
chứng trong bộ dữ liệu đánh giá, tại các mức cắt $K \in \{3, 5, 7\}$.

\begin{table}[H]
\centering
\small
\begin{tabular}{|p{2.6cm}|p{4.8cm}|p{6.3cm}|}
\hline
\textbf{Chỉ số} & \textbf{Ý nghĩa} & \textbf{Vai trò} \\
\hline
hit@K & Có ít nhất một đoạn đúng trong top-K & Đo khả năng ``tìm thấy'' \\
\hline
precision@K & Tỷ lệ đoạn đúng trong K vị trí & Đo độ ``sạch'' của ngữ cảnh \\
\hline
recall@K & Tỷ lệ bằng chứng được phủ trong top-K & Chỉ số bao phủ then chốt cho RAG \\
\hline
MRR@K & Nghịch đảo thứ hạng đoạn đúng đầu tiên & Đo ``đoạn đúng có nằm trên cao không'' \\
\hline
nDCG@K & Vừa đúng vừa đúng thứ tự, có chiết khấu logarit & Chỉ số xếp hạng tổng hợp tốt nhất \\
\hline
\end{tabular}
\caption{Định nghĩa các chỉ số truy hồi}
\label{tab:ch6-retrieval-metrics}
\end{table}

\subsection{Chỉ số sinh câu trả lời}

Các chỉ số đánh giá đầu cuối được gán bởi mô hình đánh giá tự động và bộ so khớp đáp án vàng.

\begin{table}[H]
\centering
\small
\begin{tabular}{|p{3.6cm}|p{10.1cm}|}
\hline
\textbf{Chỉ số} & \textbf{Ý nghĩa} \\
\hline
Độ bám nguồn (Faithfulness) & Câu trả lời có bám sát ngữ cảnh truy hồi không \\
\hline
Tỷ lệ ảo giác (Hallucination rate) & Tỷ lệ câu trả lời chứa nội dung không có trong nguồn --- phần bù của độ bám nguồn, càng thấp càng tốt \\
\hline
Độ liên quan (Relevance) & Câu trả lời có đúng trọng tâm câu hỏi không \\
\hline
Độ đầy đủ (Completeness) & Có trả lời đầy đủ các ý cần thiết không \\
\hline
Độ đúng so với đáp án vàng (Correctness) & So với đáp án mẫu: đúng / đúng một phần / sai --- chỉ số sát người dùng nhất \\
\hline
\end{tabular}
\caption{Định nghĩa các chỉ số sinh câu trả lời}
\label{tab:ch6-e2e-metrics}
\end{table}

% =========================================================
\section{Kết quả thực nghiệm}
% =========================================================

\subsection{Kết quả tổng thể giữa các cấu hình hợp nhất}

Bảng~\ref{tab:ch6-main} trình bày kết quả tổng hợp theo micro-average trên
15 bộ dữ liệu đang được dùng trong báo cáo. Giá trị in đậm là cấu hình tốt
hơn ở mỗi chỉ số. Đối với tỷ lệ ảo giác, giá trị thấp hơn là tốt hơn.

\begin{table}[H]
\centering
\small
\begin{tabular}{|p{2.2cm}|p{2.7cm}|>{\centering\arraybackslash}p{2.4cm}|>{\centering\arraybackslash}p{2.2cm}|>{\centering\arraybackslash}p{2.2cm}|}
\hline
\textbf{Nhóm} & \textbf{Chỉ số} & \textbf{Linear 0,8/0,2} & \textbf{RRF 0,8/0,2} & \textbf{RRF 0,5/0,5} \\
\hline
Truy hồi & nDCG@3 & \textbf{0,736} & \textbf{0,736} & 0,684 \\
& nDCG@5 & \textbf{0,757} & 0,755 & 0,712 \\
& nDCG@7 & \textbf{0,765} & 0,762 & 0,716 \\
& Recall@5 & \textbf{0,812} & 0,804 & 0,783 \\
& MRR@5 & \textbf{0,754} & 0,753 & 0,701 \\
& hit@5 & \textbf{0,843} & 0,834 & 0,811 \\
\hline
Sinh câu trả lời & Độ bám nguồn & 0,926 & \textbf{0,928} & 0,845 \\
& Tỷ lệ ảo giác & \textbf{0,043} & 0,053 & 0,136 \\
& Độ đầy đủ & 0,943 & \textbf{0,949} & 0,755 \\
& Độ liên quan & \textbf{0,996} & \textbf{0,996} & 0,847 \\
& Độ đúng & 0,845 & \textbf{0,862} & 0,819 \\
\hline
\end{tabular}
\caption{So sánh các cấu hình hợp nhất trên 15 bộ dữ liệu đánh giá (micro-average)}
\label{tab:ch6-main}
\end{table}

\subsection{Kết quả tách thành phần}

Để làm rõ đóng góp của từng thành phần trong pipeline, Bảng~\ref{tab:ch6-ablation-retrieval}
và Bảng~\ref{tab:ch6-ablation-generation} so sánh bốn thiết lập: tìm kiếm vector
thuần túy, tìm kiếm từ khóa thuần túy, RRF 0,8/0,2 nhưng tắt reranker, và RRF
0,8/0,2 đầy đủ. Các giá trị trong bảng cũng được tính theo micro-average.

\begin{table}[H]
\centering
\scriptsize
\begin{tabular}{|p{2.0cm}|>{\centering\arraybackslash}p{2.0cm}|>{\centering\arraybackslash}p{2.0cm}|>{\centering\arraybackslash}p{2.7cm}|>{\centering\arraybackslash}p{2.1cm}|}
\hline
\textbf{Chỉ số} & \textbf{Vector-only} & \textbf{Keyword-only} & \textbf{No-rerank 0,8/0,2} & \textbf{RRF 0,8/0,2} \\
\hline
nDCG@3 & 0,622 & 0,608 & 0,506 & \textbf{0,736} \\
nDCG@5 & 0,655 & 0,641 & 0,540 & \textbf{0,755} \\
nDCG@7 & 0,664 & 0,649 & 0,558 & \textbf{0,762} \\
Recall@5 & 0,719 & 0,699 & 0,651 & \textbf{0,804} \\
Hit@5 & 0,745 & 0,732 & 0,683 & \textbf{0,834} \\
MRR@5 & 0,644 & 0,634 & 0,516 & \textbf{0,753} \\
\hline
\end{tabular}
\caption{So sánh các chỉ số truy hồi trong thí nghiệm tách thành phần}
\label{tab:ch6-ablation-retrieval}
\end{table}

\begin{table}[H]
\centering
\scriptsize
\begin{tabular}{|p{3.0cm}|>{\centering\arraybackslash}p{1.8cm}|>{\centering\arraybackslash}p{1.8cm}|>{\centering\arraybackslash}p{2.6cm}|>{\centering\arraybackslash}p{2.0cm}|}
\hline
\textbf{Chỉ số} & \textbf{Vector-only} & \textbf{Keyword-only} & \textbf{No-rerank 0,8/0,2} & \textbf{RRF 0,8/0,2} \\
\hline
Độ đúng (Correctness) & 72,8\% & 62,8\% & 68,7\% & \textbf{86,2\%} \\
Độ bám nguồn (Faithfulness) & 87,7\% & 71,5\% & 77,2\% & \textbf{92,8\%} \\
Tỷ lệ ảo giác & 10,6\% & 26,2\% & 19,1\% & \textbf{5,3\%} \\
Độ đầy đủ (Completeness) & 87,4\% & 72,3\% & 78,9\% & \textbf{94,9\%} \\
\hline
\end{tabular}
\caption{So sánh các chỉ số sinh câu trả lời trong thí nghiệm tách thành phần}
\label{tab:ch6-ablation-generation}
\end{table}

% =========================================================
\section{Phân tích kết quả}
% =========================================================

\subsection{Vai trò của tìm kiếm vector và tìm kiếm từ khóa}

Kết quả tách thành phần cho thấy cả hai phương pháp đơn lẻ đều chưa đủ mạnh
khi đứng một mình. Vector-only đạt nDCG@5 là 0,655 và Recall@5 là 0,719, nhỉnh
hơn keyword-only với nDCG@5 là 0,641 và Recall@5 là 0,699. Điều này cho thấy
mô hình nhúng có lợi thế trong các câu hỏi diễn đạt tự nhiên hoặc có cách hỏi
khác với văn bản gốc. Ngược lại, tìm kiếm từ khóa vẫn có vai trò quan trọng đối
với các truy vấn chứa mã học phần, số hiệu văn bản, mốc thời gian hoặc các cụm
từ hành chính cần khớp chính xác.

Việc một phương pháp đơn lẻ không đạt kết quả tốt nhất là phù hợp với đặc thù
của dữ liệu học vụ. Nhiều câu hỏi cần đồng thời hai loại tín hiệu: khớp chính
xác các định danh như mã học phần, học kỳ, số quyết định; và hiểu quan hệ ngữ
nghĩa giữa câu hỏi của sinh viên với cách diễn đạt trong văn bản quy định. Vì
vậy, kiến trúc truy hồi lai vẫn là lựa chọn hợp lý hơn so với chỉ dùng một
nguồn tín hiệu.

\subsection{Vai trò của reranker trong pipeline}

Khi tắt reranker, chất lượng truy hồi giảm rõ rệt: nDCG@5 chỉ đạt 0,540 và
Recall@5 đạt 0,651, thấp hơn đáng kể so với RRF 0,8/0,2 đầy đủ (0,755 và
0,804). Chất lượng câu trả lời cũng bị ảnh hưởng trực tiếp, với độ đúng giảm
từ 86,2\% xuống 68,7\%. Kết quả này cho thấy reranker không chỉ là bước tinh
chỉnh thứ tự tài liệu, mà còn là thành phần quyết định trong việc đưa bằng
chứng phù hợp lên các vị trí đầu, nơi mô hình sinh câu trả lời sử dụng nhiều
nhất.

\subsection{So sánh Linear và RRF ở trọng số 0,8/0,2}

Ở cùng trọng số vector/keyword 0,8/0,2, Linear và RRF cho kết quả khá sát nhau
ở tầng truy hồi. Linear nhỉnh nhẹ ở nDCG@5 (0,757 so với 0,755), Recall@5
(0,812 so với 0,804) và tỷ lệ ảo giác thấp hơn (4,3\% so với 5,3\%). Ngược
lại, RRF đạt độ đúng so với đáp án vàng cao hơn, 86,2\% so với 84,5\%, đồng
thời có độ bám nguồn và độ đầy đủ cao hơn một chút.

Kết quả này không cho phép kết luận RRF vượt trội tuyệt đối trên mọi tiêu chí.
Thay vào đó, có thể xem RRF 0,8/0,2 là cấu hình cân bằng tốt hơn khi ưu tiên
độ đúng cuối cùng của câu trả lời, còn Linear 0,8/0,2 là một đối chứng mạnh ở
các chỉ số truy hồi thuần túy và mức ảo giác.

\subsection{Ảnh hưởng của trọng số cân bằng 0,5/0,5}

Cấu hình RRF 0,5/0,5 được đánh giá nhằm kiểm tra giả thuyết rằng tăng tỷ trọng
BM25 có thể cải thiện các câu hỏi chứa định danh chính xác. Tuy nhiên, trên 15
bộ dữ liệu hiện tại, cấu hình này thấp hơn RRF 0,8/0,2 ở hầu hết các chỉ số:
nDCG@5 đạt 0,712 so với 0,755, Recall@5 đạt 0,783 so với 0,804, và độ đúng đạt
81,9\% so với 86,2\%. Tỷ lệ ảo giác cũng tăng từ 5,3\% lên 13,6\%.

Dù vậy, RRF 0,5/0,5 vẫn vượt các cấu hình đơn lẻ về truy hồi và độ đúng: nDCG@5
0,712 cao hơn vector-only (0,655) và keyword-only (0,641), còn độ đúng 81,9\%
cũng cao hơn 72,8\% và 62,8\%. Do đó, kết quả này không phủ nhận giá trị của
truy hồi lai; nó cho thấy với tập dữ liệu hiện tại, tăng keyword weight lên 0,5
chưa phải lựa chọn tối ưu.

\subsection{Nút thắt nằm ở truy hồi và xếp hạng}

Một quan sát đáng chú ý là khi cấu hình truy hồi tốt, độ liên quan của câu trả
lời duy trì ở mức rất cao (99,6\% với cả Linear 0,8/0,2 và RRF 0,8/0,2), cho
thấy mô hình sinh có khả năng bám khá tốt vào ngữ cảnh được cung cấp. Trong khi
đó, các thay đổi ở tầng truy hồi và xếp hạng lại làm nDCG@5 dao động từ 0,540
(no-rerank) đến 0,755--0,757 (các cấu hình đầy đủ). Vì vậy, hướng cải thiện
quan trọng tiếp theo nên tập trung vào chất lượng phân đoạn, bao phủ tài liệu,
định tuyến miền và xếp hạng lại, thay vì chỉ tinh chỉnh prompt sinh câu trả
lời.

\subsection{Khuyến nghị cấu hình}

Dựa trên kết quả hiện tại, cấu hình RRF 0,8/0,2 có reranker là lựa chọn phù hợp
nhất nếu ưu tiên độ đúng của câu trả lời cuối cùng. Linear 0,8/0,2 vẫn là một
đối chứng mạnh và có tỷ lệ ảo giác thấp hơn, nhưng RRF 0,8/0,2 đạt correctness
cao nhất trong nhóm fusion. Cấu hình RRF 0,5/0,5 nên được xem là thí nghiệm độ
nhạy trọng số, chưa nên thay thế cấu hình vận hành mặc định.

% =========================================================
\section{Hạn chế}
% =========================================================

\textbf{Phạm vi số liệu tạm thời.} Báo cáo hiện sử dụng 15/16 bộ dữ liệu đánh
giá do bộ \texttt{rag\_evaluation\_dataset} đang được đánh giá lại ở cấu hình
RRF 0,5/0,5. Khi kết quả của bộ này hoàn tất, các bảng tổng hợp cần được tính
lại trên đủ 16 bộ dữ liệu để có kết luận cuối cùng. Tuy nhiên, trong phạm vi
hiện tại, tất cả cấu hình được so sánh đều sử dụng cùng một tập 15 bộ dữ liệu,
do đó tính công bằng giữa các cấu hình vẫn được bảo đảm.

\textbf{Giới hạn của bộ dữ liệu đánh giá.} Đáp án chuẩn (ground-truth) của
mỗi câu hỏi hiện được sinh từ một tệp đoạn duy nhất, chưa tính đến tính đa
nguồn và chéo kho tri thức của thông tin học vụ. Một câu hỏi về điều kiện tốt
nghệp thực tế có thể được trả lời đúng từ nhiều đoạn thuộc nhiều tệp khác
nhau, nhưng bộ đánh giá chỉ ghi nhận đoạn từ một nguồn duy nhất là kết quả
đúng. Điều này khiến các chỉ số truy hồi bị đánh giá thấp hơn năng lực thực
tế. Hướng khắc phục là xây dựng đáp án chuẩn đa nguồn --- bằng chứng cho mỗi
câu hỏi đến từ nhiều tệp đoạn và nhiều kho tri thức --- để các chỉ số phản ánh
chính xác hơn.

\textbf{Ảo giác vẫn phụ thuộc mạnh vào cấu hình truy hồi.} Với các cấu hình
tốt nhất, tỷ lệ ảo giác nằm ở mức tương đối thấp (4,3\% với Linear 0,8/0,2 và
5,3\% với RRF 0,8/0,2). Tuy nhiên, khi thay đổi trọng số sang RRF 0,5/0,5 hoặc
tắt reranker, tỷ lệ này tăng lên đáng kể. Điều này cho thấy việc kiểm soát ảo
giác không chỉ phụ thuộc vào prompt sinh câu trả lời, mà còn phụ thuộc trực
tiếp vào chất lượng ngữ cảnh được truy hồi và thứ tự các đoạn bằng chứng.

\textbf{Phạm vi thực nghiệm chưa bao phủ toàn bộ tính năng.} Kết quả trong
chương này chủ yếu là của luồng xử lý RAG cốt lõi. Khả năng tìm kiếm web dự phòng, tác nhân lập kế hoạch, và tính năng phát trực tuyến (streaming) chưa được đo lường định lượng một cách toàn diện trong bộ thực nghiệm hiện tại do yêu cầu thiết kế đáp án chuẩn phức tạp hơn cho các tác vụ đa bước.

\textbf{Không đo hiệu năng thời gian phản hồi.} Thực nghiệm được chạy trong
môi trường phát triển cục bộ, với bộ xếp hạng lại và mô hình ngôn ngữ chạy
trên CPU thay vì thiết bị tăng tốc. Số liệu độ trễ đo được trong điều kiện này
không phản ánh hiệu năng của hệ thống khi triển khai thực tế, do đó không được
đưa vào so sánh chính thức. Đánh giá độ trễ ở môi trường vận hành là hướng công
việc cần bổ sung khi hệ thống được đưa ra ngoài.

% =========================================================
\section{Kết luận chương}
% =========================================================

\subsection{Kết luận}

Thực nghiệm đầu cuối trên 15 bộ dữ liệu đánh giá cho thấy hệ thống trợ lý học
vụ đạt kết quả khả quan ở cả tầng truy hồi lẫn sinh câu trả lời. Trong nhóm
cấu hình hợp nhất đầy đủ, RRF 0,8/0,2 đạt độ đúng cao nhất so với đáp án vàng
(86,2\%), đồng thời duy trì độ bám nguồn 92,8\% và độ đầy đủ 94,9\%. Linear
0,8/0,2 nhỉnh nhẹ ở một số chỉ số truy hồi và có tỷ lệ ảo giác thấp hơn, vì
vậy hai cấu hình này cần được nhìn nhận như hai lựa chọn có đánh đổi, không
phải quan hệ thắng tuyệt đối trên mọi tiêu chí.

Kết quả tách thành phần khẳng định vai trò quan trọng của kiến trúc truy hồi
lai và reranker. Vector-only và keyword-only đều thấp hơn rõ rệt so với cấu
hình đầy đủ, trong khi việc tắt reranker làm nDCG@5 giảm xuống 0,540 và độ đúng
giảm xuống 68,7\%. Thí nghiệm RRF 0,5/0,5 cho thấy tăng trọng số từ khóa lên
mức cân bằng không cải thiện kết quả tổng thể trên tập dữ liệu hiện tại; cấu
hình này vẫn tốt hơn các phương pháp đơn lẻ, nhưng thấp hơn RRF 0,8/0,2 ở cả
truy hồi, độ đúng và tỷ lệ ảo giác. Do đó, cấu hình được khuyến nghị ở thời
điểm hiện tại là RRF 0,8/0,2 có reranker, đồng thời tiếp tục hoàn thiện dữ liệu
và đánh giá lại khi đủ 16 bộ dữ liệu.

\end{document}
